\documentclass[12pt,a4paper]{article}
\usepackage[utf8]{inputenc}
\usepackage[T5]{fontenc}
\usepackage[vietnamese]{babel}
\usepackage[margin=1.3cm,top=1.0cm,bottom=1.05cm,includehead,headheight=18pt,headsep=8pt]{geometry}
\usepackage{amsmath,amssymb,array,tabularx,enumitem,fancyhdr,tikz,xcolor,needspace}
\usetikzlibrary{calc,angles,quotes,arrows.meta,babel,patterns}
\definecolor{ink}{RGB}{30,55,70}

\newcommand{\headerQuote}{%
  \ifcase\value{page}\relax
    \or Math gives us hope that every problem has a solution% Trang 1
    \or Math + You = Champion% Trang 2
    \or Where Frustration Becomes Inspiration% Trang 3
    \or Turning Confusion into Clarity% Trang 4
    \or Don't be an X, be a Y of mathematics% Trang 5
    \or Don't get mad. Learn to add% Trang 6
    \or The only way to learn mathematics is to do mathematics% Trang 7
    \or Small steps every day lead to big results% Trang 8
    \or Every expert was once a beginner% Trang 9
    \or Success is the sum of small efforts repeated daily% Trang 10
    \or Mathematics is the music of reason% Trang 11
    \or Believe you can and you're halfway there% Trang 12
    \else Math gives us hope that every problem has a solution%
  \fi
}
\pagestyle{fancy}\fancyhf{}
\lhead{\small\textit{\headerQuote}}
\rhead{\small\textit{Người biên soạn: Dương Gia Bảo}}
\rfoot{\small Trang \thepage}
\renewcommand{\headrulewidth}{0.4pt}
\setlength{\headheight}{18pt}
\setlength{\parindent}{0pt}
\setlength{\parskip}{0pt}
\setlength{\tabcolsep}{1.5pt}
\renewcommand{\arraystretch}{1.0}
\newcommand{\parttitle}[1]{\par\vspace{1.5pt}\needspace{2\baselineskip}{\large\bfseries\color{ink}#1}\par\nopagebreak\vspace{1pt}}
\newcommand{\four}[4]{\par\noindent\begin{tabularx}{\linewidth}{@{}*{4}{>{\raggedright\arraybackslash}X}@{}}\textbf{A.} #1&\textbf{B.} #2&\textbf{C.} #3&\textbf{D.} #4\end{tabularx}}
\newcommand{\two}[4]{\par\noindent\begin{tabularx}{\linewidth}{@{}*{2}{>{\raggedright\arraybackslash}X}@{}}\textbf{A.} #1&\textbf{B.} #2\\[1pt]\textbf{C.} #3&\textbf{D.} #4\end{tabularx}}
\newcommand{\answerbox}{\hfill\textit{Đáp số:}\ \fbox{\rule{0pt}{9pt}\hspace{2.2cm}}}
\setlist[enumerate,1]{label=\textbf{Câu \arabic*.},leftmargin=1.1cm,labelsep=.1cm,itemsep=0.2pt,topsep=0.2pt,parsep=0pt}
\setlist[enumerate,2]{label=\alph*),leftmargin=.55cm,itemsep=0.2pt,topsep=0.2pt,parsep=0pt}

\newcommand{\soDe}{N}

\begin{document}
\noindent
\begin{minipage}[c]{0.38\textwidth}
    \centering
    \textbf{LÀM TOÁN ĐỂ HẠNH PHÚC}\\[-2pt]
    \rule{3.2cm}{0.6pt}\\[2pt]
    \textbf{ĐỀ CHÍNH THỨC}
\end{minipage}\hfill
\begin{minipage}[c]{0.35\textwidth}
    \centering
    \textbf{ĐỀ ÔN THI GIỮA KÌ I}\\
    \textbf{NĂM HỌC 2026 -- 2027}\\[1pt]
    \textbf{ĐỀ SỐ \soDe}\\
    \textbf{Môn: TOÁN -- Khối lớp 10}\\
    \textit{\small Thời gian làm bài: 90 phút\\(không kể thời gian phát đề)}
\end{minipage}

\vspace{0.04cm}
\noindent
\begin{minipage}[b]{0.78\textwidth}
    \textbf{Họ và tên:}\dotfill\textbf{SBD:}\makebox[3.5cm]{\dotfill}
\end{minipage}\hfill
\begin{minipage}[b]{0.20\textwidth}
    \hfill
    \setlength{\fboxrule}{0.8pt}%
    \fbox{\begin{minipage}{2.3cm}
        \centering
        \textbf{\small Mã đề thi}\\
        \textbf{00\soDe}
    \end{minipage}}
\end{minipage}
\vspace{0.04cm}

% =========================================================================
% TRANG 1: PHẦN I (12 CÂU TRẮC NGHIỆM NHIỀU LỰA CHỌN)
% Quy chuẩn: 12 câu nằm gọn 100% trong Trang 1, kết thúc bằng \newpage.
% Cân bằng đáp án: 3A - 3B - 3C - 3D.
% =========================================================================
\parttitle{PHẦN I. Trắc nghiệm nhiều lựa chọn}
\textit{Thí sinh trả lời từ câu 1 đến câu 12. Mỗi câu hỏi chỉ chọn một phương án đúng duy nhất.}
\begin{enumerate}
% [Điền 12 câu trắc nghiệm nhiều lựa chọn tại đây]
\item Câu 1...
\four{A.}{B.}{C.}{D.}
% ...
\item Câu 12...
\two{A.}{B.}{C.}{D.}
\end{enumerate}
\newpage

% =========================================================================
% TRANG 2: PHẦN II (2 CÂU ĐÚNG/SAI) + PHẦN III (CÂU 1 & CÂU 2 CÓ HÌNH TIKZ)
% Quy chuẩn: Nằm trọn vẹn trong Trang 2, kết thúc bằng \newpage.
% Hình TikZ đặt căn giữa dưới đề bài (\begin{center}\begin{tikzpicture}...\end{tikzpicture}\end{center}).
% Chống đè chữ: Trục Ox số đo below, đỉnh above; Trục Oy số đo left, đỉnh above right;
% Nhãn đường thẳng đặt ở đầu mút xa thoáng (above left/above right).
% =========================================================================
\parttitle{PHẦN II. Trắc nghiệm đúng/sai}
\textit{Thí sinh trả lời từ câu 1 đến câu 2. Trong mỗi ý a), b), c), d) ở mỗi câu, thí sinh chọn đúng hoặc sai.}
\begin{enumerate}
\item Câu 1 đúng/sai...
\begin{enumerate}
    \item Ý a...
    \item Ý b...
    \item Ý c...
    \item Ý d...
\end{enumerate}
\item Câu 2 đúng/sai...
\begin{enumerate}
    \item Ý a...
    \item Ý b...
    \item Ý c...
    \item Ý d...
\end{enumerate}
\end{enumerate}

\parttitle{PHẦN III. Trắc nghiệm trả lời ngắn}
\textit{Thí sinh trả lời từ câu 1 đến câu 4. Điền kết quả vào ô tương ứng.}
\begin{enumerate}
\item Câu 1 trả lời ngắn (bài toán có hình vẽ TikZ căn giữa bên dưới)...
\begin{center}
\vspace{-3pt}
% \begin{tikzpicture}[scale=0.8, >=Stealth] ... \end{tikzpicture}
\vspace{-6pt}
\end{center}
\answerbox

\item Câu 2 trả lời ngắn (bài toán nửa mặt phẳng bờ d hoặc hình vẽ TikZ)...
\begin{center}
\vspace{-3pt}
% \begin{tikzpicture}[scale=0.72, >=Stealth] ... \end{tikzpicture}
\vspace{-6pt}
\end{center}
\answerbox
\newpage

% =========================================================================
% TRANG 3: PHẦN III (CÂU 3 & CÂU 4) + PHẦN IV (4 CÂU TỰ LUẬN)
% Quy chuẩn: Nằm trọn vẹn trong Trang 3, kết thúc bằng \newpage.
% =========================================================================
\item Câu 3 trả lời ngắn...
\answerbox

\item Câu 4 trả lời ngắn...
\answerbox
\end{enumerate}

\parttitle{PHẦN IV. Tự luận}
\textit{Thí sinh trình bày chi tiết lời giải từ câu 1 đến câu 4.}
\begin{enumerate}
\item \textit{(0,75 điểm)} Câu 1 tự luận...
\item \textit{(0,75 điểm)} Câu 2 tự luận...
\item \textit{(0,75 điểm)} Câu 3 tự luận...
\item \textit{(0,75 điểm)} Câu 4 tự luận (bài toán thực tế tối ưu hóa tuyến tính)...
\end{enumerate}
\newpage

% =========================================================================
% TRANG 4 TRỞ ĐI: BẢNG ĐÁP ÁN TỔNG HỢP VÀ LỜI GIẢI CHI TIẾT
% =========================================================================
\parttitle{BẢNG ĐÁP ÁN TỔNG HỢP}

\vspace{0.1cm}
\textbf{1. Phần I: Trắc nghiệm nhiều lựa chọn}
\begin{center}
{\setlength{\tabcolsep}{2.5pt}\small
\begin{tabular}{|c|c|c|c|c|c|c|c|c|c|c|c|}
\hline
\textbf{Câu 1} & \textbf{Câu 2} & \textbf{Câu 3} & \textbf{Câu 4} & \textbf{Câu 5} & \textbf{Câu 6} & \textbf{Câu 7} & \textbf{Câu 8} & \textbf{Câu 9} & \textbf{Câu 10} & \textbf{Câu 11} & \textbf{Câu 12} \\
\hline
... & ... & ... & ... & ... & ... & ... & ... & ... & ... & ... & ... \\
\hline
\end{tabular}}
\end{center}

\vspace{0.2cm}
\textbf{2. Phần II: Trắc nghiệm đúng/sai}
\begin{center}
{\small
\begin{tabular}{|c|c|c|c|c|}
\hline
\textbf{Câu} & \textbf{Lệnh hỏi a)} & \textbf{Lệnh hỏi b)} & \textbf{Lệnh hỏi c)} & \textbf{Lệnh hỏi d)} \\
\hline
\textbf{1} & Đ/S & Đ/S & Đ/S & Đ/S \\
\hline
\textbf{2} & Đ/S & Đ/S & Đ/S & Đ/S \\
\hline
\end{tabular}}
\end{center}

\vspace{0.2cm}
\textbf{3. Phần III: Trắc nghiệm trả lời ngắn}
\begin{center}
{\small
\begin{tabular}{|c|c|c|c|c|}
\hline
\textbf{Câu} & \textbf{Câu 1} & \textbf{Câu 2} & \textbf{Câu 3} & \textbf{Câu 4} \\
\hline
\textbf{Đáp số} & ... & ... & ... & ... \\
\hline
\end{tabular}}
\end{center}

\vspace{0.2cm}
\parttitle{LỜI GIẢI CHI TIẾT}
% [Điền lời giải chi tiết từng câu tại đây]
% Với bài toán hệ BPT / tối ưu thực tế (Câu 4 Tự luận): Bắt buộc vẽ đồ thị miền nghiệm TikZ to rõ,
% áp dụng triệt để quy tắc chống đè chữ (Anti-collision TikZ Rule).

\end{document}
